\documentclass[11pt]{article}
\usepackage[a4paper,margin=22mm]{geometry}
\usepackage{fontspec}
\setmainfont{DejaVu Serif}
\setsansfont{DejaVu Sans}
\usepackage{microtype}
\usepackage{xcolor}
\usepackage{hyperref}
\usepackage{enumitem}
\usepackage{parskip}
\definecolor{Accent}{HTML}{971D32}
\definecolor{Muted}{HTML}{666666}
\hypersetup{colorlinks=true,urlcolor=Accent,linkcolor=Accent}
\setlist{leftmargin=1.6em,itemsep=0.3em,topsep=0.35em}
\setcounter{tocdepth}{2}
\renewcommand{\contentsname}{Contents}
\newcommand{\sectionrule}{\par\vspace{-0.5em}\noindent\textcolor{Accent}{\rule{\linewidth}{0.6pt}}\par}
\begin{document}

\begin{center}
{\sffamily\bfseries\color{Accent} TOEFL WORD OF THE DAY\par}
\vspace{8mm}
{\fontsize{42}{48}\selectfont\bfseries concurrent\par}
\vspace{2mm}
{\Large\sffamily\color{Accent} /kənˈkɝːənt/\par}
\vspace{2mm}
{\small\sffamily Local audio phrase: concurrent\par}
{\small\sffamily Audio file: audios/concurrent.mp3\par}
\vspace{2mm}
{\large\sffamily\color{Muted} adjective\par}
\vspace{5mm}
{\sffamily happening at the same time \textbullet\ shared by multiple authorities\par}
\vspace{6mm}
{\large Concurrent describes events that happen together or powers and responsibilities held by more than one authority at the same time.\par}
\vspace{6mm}
\begin{minipage}{0.84\textwidth}
\itshape “The university scheduled three concurrent workshops, so students had to choose one.”
\end{minipage}\par
\vspace{10mm}
\textbf{Date:} July 31st 2026\quad\textbullet\quad \textbf{Level:} B1--mid-B2 TOEFL
\vfill
Created and compiled by \textbf{Amir Kooshky}\par
\vspace{2mm}
\href{https://t.me/KooshkyTOEFL}{Telegram Channel}\quad\textbullet\quad
\href{https://www.instagram.com/kooshkytoefl}{Instagram}\quad\textbullet\quad
\href{https://t.me/KooshkyTOEFL\_pv}{Personal Telegram}
\end{center}
\newpage
\tableofcontents
\newpage

\section{Meanings and usage}
\sectionrule
\subsection{Meaning 1: Happening at the same time}
\textbf{Part of speech:} adjective

\textbf{IPA:} /kənˈkɝːənt/

\textbf{Pronunciation phrase:} concurrent

\textbf{Local audio file:} audios/concurrent.mp3

\textbf{Register:} formal; common in academic, medical, technological, legal, and professional contexts

\textbf{Definition:} happening, existing, or being done at the same time as something else

\textbf{In simpler words:} happening at the same time

\textbf{Typical pattern:} concurrent with + event or concurrent + events, processes, or activities

\subsubsection{Natural examples}
\begin{enumerate}
  \item The university scheduled three concurrent workshops, so students had to choose one.
  \item The patient was treated for two concurrent infections.
  \item The software can perform several concurrent tasks without slowing down.
  \item The researchers monitored concurrent changes in temperature and humidity.
  \item Two concurrent construction projects caused severe traffic delays.
  \item She completed a degree while receiving concurrent professional training.
  \item The conference offered concurrent sessions on language learning and educational technology.
\end{enumerate}

\subsubsection{Nuance and implication}
\textbf{Central nuance:} The word emphasizes that two or more things occur during the same period rather than one after another.

\textbf{Usual implication:} The events may begin and end at different moments, but at least part of their duration overlaps.

\textbf{Compared with consecutive:} Concurrent events happen at the same time, while consecutive events follow one another without occurring together.

\subsubsection{Grammar notes}
\textbf{How to use it:} Use it before a noun, as in concurrent events, or after a linking verb, as in The projects were concurrent.

\textbf{What can follow it:} It commonly appears before words such as event, session, process, task, illness, treatment, sentence, and program.

\textbf{Common preposition:} Use concurrent with when comparing one event or condition with another, as in The training was concurrent with her university studies.

\textbf{Position in a sentence:} Concurrent usually appears directly before the noun it describes.

\subsubsection{Useful patterns}
\textbf{Pattern:} concurrent + events, sessions, or processes

\textbf{Use:} Use this for activities or developments occurring during the same period.

\textbf{Example:} Delegates could choose between four concurrent sessions.

\textbf{Pattern:} be concurrent with + noun

\textbf{Use:} Use this to state that one event happens during the same period as another.

\textbf{Example:} The population decline was concurrent with a prolonged drought.

\textbf{Pattern:} run or occur concurrently

\textbf{Use:} Use the adverb concurrently when describing how two or more activities happen.

\textbf{Example:} The two experiments were conducted concurrently.

\textbf{Pattern:} concurrent sentences

\textbf{Use:} In law, this describes prison sentences that are served at the same time rather than one after another.

\textbf{Example:} The judge ordered the two prison sentences to run concurrently.

\subsubsection{Common wrong usage}
\paragraph{Correction 1}
\textbf{Wrong:} The two courses are concurrent to each other.

\textbf{Correct:} The two courses are concurrent with each other.

\textbf{Why:} Use concurrent with, not concurrent to.

\paragraph{Correction 2}
\textbf{Wrong:} The first lecture ended, and then a concurrent lecture began.

\textbf{Correct:} The first lecture ended, and then a consecutive lecture began.

\textbf{Why:} Concurrent means happening at the same time. Use consecutive when one event follows another.

\paragraph{Correction 3}
\textbf{Wrong:} The meetings happened concurrent.

\textbf{Correct:} The meetings happened concurrently.

\textbf{Why:} Use the adverb concurrently to describe how something happened.

\subsection{Meaning 2: Shared by multiple authorities}
\textbf{Part of speech:} adjective

\textbf{IPA:} /kənˈkɝːənt/

\textbf{Pronunciation phrase:} concurrent

\textbf{Local audio file:} audios/concurrent.mp3

\textbf{Register:} formal and specialized; especially common in law, government, and administration

\textbf{Definition:} held or exercised at the same time by two or more authorities, organizations, or people

\textbf{In simpler words:} belonging to more than one authority at once

\textbf{Typical pattern:} concurrent jurisdiction, powers, authority, control, or responsibility

\subsubsection{Natural examples}
\begin{enumerate}
  \item Federal and state courts may have concurrent jurisdiction over certain cases.
  \item The constitution gives the national and regional governments some concurrent powers.
  \item Both departments have concurrent responsibility for protecting the site.
  \item The agencies exercise concurrent authority over environmental regulation.
  \item The agreement established concurrent control of the shared research facility.
  \item Several institutions have concurrent roles in managing the public health system.
  \item The two committees were given concurrent responsibility for reviewing the proposal.
\end{enumerate}

\subsubsection{Nuance and implication}
\textbf{Central nuance:} The word indicates that the same type of power, control, or responsibility exists in more than one place at the same time.

\textbf{Usual implication:} The authorities may act separately within the same area rather than always acting as a single combined group.

\textbf{Compared with exclusive:} Concurrent authority is held by more than one body, while exclusive authority belongs to only one body.

\subsubsection{Grammar notes}
\textbf{How to use it:} Use it before a noun naming the shared power, legal right, or responsibility.

\textbf{What can follow it:} Common nouns after it include jurisdiction, authority, power, control, responsibility, and role.

\textbf{Common preposition:} Jurisdiction and authority are often followed by over, as in concurrent jurisdiction over the dispute.

\textbf{Position in a sentence:} In this meaning, concurrent normally appears directly before the noun it describes.

\subsubsection{Useful patterns}
\textbf{Pattern:} concurrent jurisdiction over + matter

\textbf{Use:} Use this when two or more courts or governments have legal power over the same type of case.

\textbf{Example:} The two courts have concurrent jurisdiction over the dispute.

\textbf{Pattern:} concurrent powers

\textbf{Use:} Use this for powers that more than one level of government may exercise.

\textbf{Example:} Collecting taxes is a concurrent power in some federal systems.

\textbf{Pattern:} concurrent responsibility for + noun or -ing form

\textbf{Use:} Use this when two or more groups share responsibility for something.

\textbf{Example:} The ministries have concurrent responsibility for managing the program.

\textbf{Pattern:} concurrent authority over + area

\textbf{Use:} Use this when several bodies possess authority in the same area.

\textbf{Example:} The agencies exercise concurrent authority over workplace safety.

\subsubsection{Common wrong usage}
\paragraph{Correction 1}
\textbf{Wrong:} The courts have a concurrent jurisdiction over the case.

\textbf{Correct:} The courts have concurrent jurisdiction over the case.

\textbf{Why:} Jurisdiction is normally used without a in this expression.

\paragraph{Correction 2}
\textbf{Wrong:} Both agencies have concurrent authorities over the industry.

\textbf{Correct:} Both agencies have concurrent authority over the industry.

\textbf{Why:} Authority is normally uncountable when it means official power.

\paragraph{Correction 3}
\textbf{Wrong:} Only one department has concurrent control of the system.

\textbf{Correct:} Two departments have concurrent control of the system.

\textbf{Why:} Concurrent control requires control to be held by more than one party.

\section{Word family}
\sectionrule
\subsection{concurrently}
\textbf{Part of speech:} adverb

\textbf{IPA:} /kənˈkɝːəntli/

\textbf{Definition:} at the same time

\textbf{Relationship to the headword:} This adverb describes events or actions that happen at the same time.

\textbf{Grammar pattern:} occur, run, operate, study, or work concurrently

\subsubsection{Examples}
\begin{enumerate}
  \item The two studies were conducted concurrently.
  \item Students may enroll in both programs concurrently.
\end{enumerate}

\subsection{concurrency}
\textbf{Part of speech:} uncountable noun

\textbf{IPA:} /kənˈkɝːənsi/

\textbf{Definition:} the state of two or more events, processes, or tasks happening during the same period

\textbf{Relationship to the headword:} This noun names the condition of being concurrent.

\textbf{Grammar pattern:} support, manage, or increase concurrency

\textbf{Usage note:} It is especially common in computing and other technical fields.

\subsubsection{Examples}
\begin{enumerate}
  \item The system was designed to support a high level of concurrency.
  \item Concurrency allows several computing tasks to make progress during the same period.
\end{enumerate}

\section{Common collocations}
\sectionrule
\subsection{concurrent events}
\textbf{Applies to:} Happening at the same time

\textbf{Meaning:} events occurring during the same period

\textbf{Example:} The study examined the effects of several concurrent events.

\subsection{concurrent sessions}
\textbf{Applies to:} Happening at the same time

\textbf{Meaning:} sessions scheduled for the same time

\textbf{Example:} The conference offered six concurrent sessions in the afternoon.

\subsection{concurrent processes}
\textbf{Applies to:} Happening at the same time

\textbf{Meaning:} processes operating during the same period

\textbf{Example:} The computer manages multiple concurrent processes efficiently.

\subsection{concurrent treatment}
\textbf{Applies to:} Happening at the same time

\textbf{Meaning:} two or more treatments provided during the same period

\textbf{Example:} The patient received concurrent treatment for anxiety and chronic pain.

\subsection{concurrent enrollment}
\textbf{Applies to:} Happening at the same time

\textbf{Meaning:} being enrolled in two educational programs or institutions during the same period

\textbf{Example:} Concurrent enrollment allows high school students to take university courses.

\subsection{run concurrently}
\textbf{Applies to:} Happening at the same time

\textbf{Meaning:} continue or operate at the same time

\textbf{Pattern:} two or more activities run concurrently

\textbf{Example:} The training courses will run concurrently for six weeks.

\subsection{concurrent jurisdiction}
\textbf{Applies to:} Shared by multiple authorities

\textbf{Meaning:} legal authority held by more than one court or government

\textbf{Pattern:} concurrent jurisdiction over + case or subject

\textbf{Example:} The federal and regional courts have concurrent jurisdiction over the matter.

\subsection{concurrent powers}
\textbf{Applies to:} Shared by multiple authorities

\textbf{Meaning:} powers that may be exercised by more than one authority

\textbf{Example:} The constitution identifies several concurrent powers shared by different levels of government.

\subsection{concurrent responsibility}
\textbf{Applies to:} Shared by multiple authorities

\textbf{Meaning:} responsibility held by more than one person or organization

\textbf{Pattern:} concurrent responsibility for + noun or -ing form

\textbf{Example:} The two agencies have concurrent responsibility for monitoring water quality.

\subsection{concurrent authority}
\textbf{Applies to:} Shared by multiple authorities

\textbf{Meaning:} official power held by more than one body

\textbf{Pattern:} concurrent authority over + area

\textbf{Example:} Local and national officials exercise concurrent authority over the protected area.

\section{Synonyms and antonyms}
\sectionrule
\subsection{Close synonyms}
\subsubsection{simultaneous}
\textbf{Part of speech:} adjective

\textbf{Applies to:} Happening at the same time

\textbf{Shared meaning:} Both words describe events occurring at the same time.

\textbf{Difference:} Simultaneous often emphasizes that events occur at exactly the same moment. Concurrent can describe events that overlap during the same general period, even if their starting and ending times differ.

\textbf{Example:} The two explosions were almost simultaneous.

\subsubsection{overlapping}
\textbf{Part of speech:} adjective

\textbf{Applies to:} Happening at the same time

\textbf{Shared meaning:} Both can describe events whose time periods occur together.

\textbf{Difference:} Overlapping emphasizes that only part of the time period is shared. Concurrent more generally means occurring during the same period.

\textbf{Example:} The students had two overlapping classes on Monday.

\subsubsection{parallel}
\textbf{Part of speech:} adjective

\textbf{Applies to:} Happening at the same time

\textbf{Shared meaning:} Both can describe activities proceeding at the same time.

\textbf{Difference:} Parallel processes develop or operate beside each other and may be similar or connected. Concurrent only requires them to occur during the same period.

\textbf{Example:} The researchers conducted parallel investigations in three countries.

\subsubsection{shared}
\textbf{Part of speech:} adjective

\textbf{Applies to:} Shared by multiple authorities

\textbf{Shared meaning:} Both can describe power or responsibility belonging to more than one party.

\textbf{Difference:} Shared is a broad everyday term. Concurrent is more formal and often means that each authority independently possesses the same power or responsibility.

\textbf{Example:} The two departments have shared responsibility for the project.

\subsection{Direct or useful antonyms}
\subsubsection{consecutive}
\textbf{Part of speech:} adjective

\textbf{Applies to:} Happening at the same time

\textbf{Opposition:} Consecutive events occur one after another, while concurrent events occur during the same period.

\textbf{Example:} The interviews were held on three consecutive days.

\subsubsection{sequential}
\textbf{Part of speech:} adjective

\textbf{Applies to:} Happening at the same time

\textbf{Opposition:} Sequential processes follow a particular order, while concurrent processes operate during the same period.

\textbf{Example:} The stages must be completed in sequential order.

\subsubsection{exclusive}
\textbf{Part of speech:} adjective

\textbf{Applies to:} Shared by multiple authorities

\textbf{Opposition:} Exclusive authority belongs to only one party, while concurrent authority is held by more than one.

\textbf{Example:} The national court has exclusive jurisdiction over these cases.

\section{Words learners may confuse}
\sectionrule
\subsection{consecutive}
\textbf{Part of speech:} adjective

\textbf{Why learners confuse them:} Both words are formal adjectives often used to describe the timing or order of events.

\textbf{Key difference:} Concurrent means happening at the same time, while consecutive means happening one after another.

\textbf{concurrent:} The festival held several concurrent performances.

\textbf{consecutive:} The theater presented performances on five consecutive evenings.

\section{Etymology}
\sectionrule
\textbf{Origin:} Concurrent ultimately comes from Latin concurrere, meaning to run together.

\section{Practice exercises}
\sectionrule
\subsection{Word family choice}
Choose the best member of the word family for each blank.

\begin{enumerate}
  \item The two research teams worked \_\_\_\_\_ on different parts of the project.
  \item The software supports \_\_\_\_\_, allowing several tasks to make progress during the same period.
  \item The conference included several \_\_\_\_\_ workshops.
\end{enumerate}

\subsection{Correct or incorrect?}
Decide whether each sentence is correct. Correct the inaccurate ones.

\begin{enumerate}
  \item The two courses run concurrently, so students cannot attend both.
  \item The second stage began concurrently after the first stage ended.
  \item The regional and national courts have concurrent jurisdiction over the case.
  \item The training program was concurrent to her university studies.
\end{enumerate}

\subsection{Collocation practice}
Complete each sentence with a natural collocation.

\begin{enumerate}
  \item The conference offered five \_\_\_\_\_ sessions, each covering a different subject.
  \item The two experiments were designed to run \_\_\_\_\_.
  \item Both courts have concurrent \_\_\_\_\_ over disputes involving this law.
\end{enumerate}

\subsection{Complete answer key}
\subsubsection{Word family choice}
\begin{enumerate}
  \item \textbf{concurrently.} The two research teams worked concurrently on different parts of the project. \emph{Use the adverb concurrently to describe how the teams worked.}
  \item \textbf{concurrency.} The software supports concurrency, allowing several tasks to make progress during the same period. \emph{Use the noun concurrency after supports.}
  \item \textbf{concurrent.} The conference included several concurrent workshops. \emph{Use the adjective concurrent before workshops.}
\end{enumerate}

\subsubsection{Correct or incorrect?}
\begin{enumerate}
  \item \textbf{Correct} \emph{Run concurrently correctly means that the courses take place during the same period.}
  \item \textbf{Incorrect}. The second stage began immediately after the first stage ended. \emph{Concurrent activities overlap in time; they do not occur one after another.}
  \item \textbf{Correct} \emph{Concurrent jurisdiction means that both courts possess legal authority over the case.}
  \item \textbf{Incorrect}. The training program was concurrent with her university studies. \emph{Use concurrent with, not concurrent to.}
\end{enumerate}

\subsubsection{Collocation practice}
\begin{enumerate}
  \item \textbf{concurrent.} The conference offered five concurrent sessions, each covering a different subject. \emph{Concurrent sessions take place at the same time.}
  \item \textbf{concurrently.} The two experiments were designed to run concurrently. \emph{Run concurrently is a common expression for activities operating at the same time.}
  \item \textbf{jurisdiction.} Both courts have concurrent jurisdiction over disputes involving this law. \emph{Concurrent jurisdiction is legal authority held by more than one court.}
\end{enumerate}

\vfill
\begin{center}
\small\color{Muted} Created and compiled by \textbf{Amir Kooshky}\\
\href{https://t.me/KooshkyTOEFL}{Telegram Channel}\quad\textbullet\quad
\href{https://www.instagram.com/kooshkytoefl}{Instagram}\quad\textbullet\quad
\href{https://t.me/KooshkyTOEFL\_pv}{Personal Telegram}
\end{center}
\end{document}
